\subsection{Shared experience banks and dependency-aware skill retrieval narrow the residual further}

A closer parent for the ``learn once, reuse across tasks'' intuition is ReX, which treats recurring reasoning patterns and skills as latent experiences rather than task-specific adapters \cite{ling2026rex}. ReX learns a shared Experience Bank of foundational skill vectors, encodes each input into a low-dimensional experience code, and dynamically composes the relevant skill vectors into a lightweight adapter without requiring explicit task identifiers. Its multi-task results show that cross-input/task experience reuse can improve specialization and generalization. This removes any novelty claim that Paper 3 uniquely proposes a persistent shared experience substrate or dynamic composition of reusable skills.

Dependency-aware skill retrieval is also moving beyond flat semantic matching. SkillGraph represents agent skills as a dependency graph and targets selection/composition failures that arise when modular skills are treated as isolated text descriptions \cite{wu2026skillgraph}. SKILLGRAPH goes further by representing reusable skills as an evolving typed graph whose prerequisite, enhancement and co-occurrence edges are updated from trajectories and reinforcement-learning feedback; ordered skill subgraphs guide new-task execution while the same graph participates in improving the policy \cite{li2026skillgraphrl}. GraSP compiles retrieved skills into executable typed DAGs with precondition--effect edges, node-level verification and locality-bounded repair \cite{xia2026grasp}. AgentGL similarly treats graph topology as an active reasoning substrate: an LLM agent uses graph-native tools at inference while a graph-conditioned curriculum is used during reinforcement learning \cite{sun2026agentgl}. These systems are important parent mechanisms because they show that relational topology can inform learning, retrieval, composition and execution.

A complementary retrieval parent exposes why the negative controls must be strong. SkillSight shows that recurring descriptive boilerplate across skill documents can inflate dense and lexical relevance scores even when those shared words are not discriminative of the required capability; calibrating that background improves retrieval without extra model training \cite{xiao2026skillsight}. Paper 3 therefore cannot treat a semantic retriever failure as evidence for structural novelty unless the semantic baseline is already calibrated against shared-background bias.

The candidate residual must therefore remain more specific. Paper 3 asks whether a \emph{scientifically scoped} structural object---with explicit QoI, context, directional role mapping, preserved and non-preserved properties, boundary conditions, evidence and negative-transfer history---can do two jobs with the same persistent identity: estimate \emph{cross-domain training-data redundancy} and license or reject inference-time transfer. ReX supplies reusable latent experience; SkillGraph/SKILLGRAPH/GraSP supply increasingly rich dependency and execution graphs; AgentGL demonstrates one graph-conditioned learning/execution loop; SkillSight supplies a stronger semantic-retrieval control. RAKL must either add value in the controlled surface-disjoint, boundary-sensitive training-selection plus inference regime or narrow its claim to a scientific-validity governance layer rather than a general reuse advance.
